\section{Viabilidad técnica y estado del prototipo}

\subsection{Arquitectura implementada}

GymHub utiliza una arquitectura web modular. La interfaz se construye con React
y TypeScript; React permite organizar la interfaz mediante componentes
reutilizables \parencite{react2026}. El backend se desarrolla con Django 5.1 y
Django REST Framework, que proporcionan modelos, autenticación, validación y
exposición de API \parencite{django51,drf315}. PostgreSQL almacena los datos en
Supabase.

Las tareas programadas se modelan con Celery y Redis. Una cola de tareas separa
el envío del trabajo de su ejecución por procesos worker
\parencite{celery2026}. En el entorno local esos procesos se orquestan con Docker
Compose; en la demostración publicada en Vercel no están desplegados como procesos
persistentes. Esta diferencia limita alertas y automatizaciones programadas en
producción.

\begin{table}[H]
\caption{Tecnologías verificadas en el repositorio}
\begin{tabularx}{\textwidth}{>{\bfseries}p{3cm}p{4cm}X}
\toprule
Capa & Tecnología & Uso en GymHub \\
\midrule
Frontend & React 18, TypeScript, Vite 5 & SPA responsive, rutas, formularios y vistas por rol. \\
Interfaz & Tailwind CSS, Recharts, Lucide & Estilos, gráficos e iconografía. \\
Estado y datos & TanStack Query, Zustand, Axios & Caché de servidor, sesión y comunicación HTTP. \\
Backend & Django 5.1, DRF 3.15 & Reglas de negocio, permisos y API REST. \\
Autenticación & JWT y cookies HTTP-only & Inicio, renovación y cierre de sesión por rol. \\
Datos & PostgreSQL en Supabase & Persistencia relacional y migraciones. \\
Procesos & Celery, Celery Beat y Redis & Tareas asíncronas, alertas y programación local. \\
Despliegue & Vercel y Docker & Demo pública y entornos reproducibles. \\
Pruebas & Pytest, Vitest y Playwright & Reglas backend, componentes y flujos de navegador. \\
\bottomrule
\end{tabularx}
\caption*{\textit{Nota}. Elaboración propia a partir de archivos de dependencias y configuración del repositorio \parencite{gymhub_repo}.}
\end{table}

\subsection{Módulos funcionales}

El MVP define tres capacidades de acceso: miembro, entrenador y personal
administrativo. El flujo demostrable abarca inicio de sesión, panel por rol,
miembros, planes, asistencia, progreso, pagos, alertas, nutrición, chat y perfil.

\begin{longtable}{>{\bfseries}p{3.1cm}p{7.3cm}p{3.8cm}}
\caption{Estado de los módulos del MVP}\label{tab:modulos}\\
\toprule
Módulo & Capacidad observable & Estado y límite \\
\midrule
\endfirsthead
\toprule
Módulo & Capacidad observable & Estado y límite \\
\midrule
\endhead
Autenticación & Inicio, cierre, renovación y acceso por rol. & Funcional; requiere operación segura de secretos y cookies. \\
Miembros & Lista, detalle, activación y asignación. & Funcional en MVP. \\
Planes & Planes, días, ejercicios y entrenamiento del día. & Funcional; requiere mayor ciclo editorial comercial. \\
Asistencia & Check-in e historial con reglas de mora. & Funcional; ingreso físico automatizado es futuro. \\
Progreso & Sesiones, registros de ejercicios y métricas. & Funcional en nivel básico. \\
Facturación & Planes, pagos, vencimientos y mora. & Funcional como registro; no es contabilidad ni pasarela bancaria. \\
Alertas & Inactividad y notificaciones. & Código funcional; programación persistente no desplegada en Vercel. \\
Nutrición & Perfil y guías generales. & Funcional como orientación informativa, no clínica. \\
Chat contextual & Respuestas con contexto de datos del usuario. & Funcional con reglas; proveedor LLM es opcional. \\
Gráficas & Indicadores de operación y progreso. & Funcional; archivos media no tienen almacenamiento persistente externo. \\
Clases & Clases e inscripciones en backend. & Dominio existente, fuera del flujo principal congelado del MVP. \\
\bottomrule
\caption*{\textit{Nota}. Auditoría del repositorio al 1 de julio de 2026.}
\end{longtable}

\subsection{Evidencia de avance}

La inspección del repositorio identificó 18 archivos de pruebas backend con 148
casos, 26 archivos de pruebas frontend y dos especificaciones de extremo a
extremo. Estas cantidades evidencian trabajo técnico, pero no equivalen a una
certificación de calidad ni sustituyen pruebas de aceptación con usuarios.

El frontend, backend y esquema de API disponen de despliegues de demostración.
Las credenciales internas no se incluyen en este documento. Para la exposición
se recomienda preparar un video y capturas anonimizadas como respaldo si falla la
conectividad.

\subsection{Seguridad y privacidad}

El prototipo usa permisos por rol, cookies HTTP-only para JWT, controles CORS y
CSRF, y configuración separada por ambiente. Django incluye mecanismos para
autenticación, autorización y seguridad web \parencite{django51}. La existencia
de estos mecanismos no elimina la necesidad de revisión, actualización de
dependencias, rotación de secretos, copias de seguridad y respuesta a incidentes.

Los datos de identidad, pagos, asistencia y progreso deben tratarse según
finalidad, consentimiento y acceso autorizado. La Ley n.º 8968 aplica a bases de
datos privadas y protege el control de las personas sobre su información
\parencite{ley8968}. Toda demostración utilizará datos ficticios o anonimizados.

\subsection{Alcance de salud y nutrición}

Las funciones de alimentación se limitan a organización y orientación general.
No realizan diagnóstico, tratamiento, prescripción clínica ni evaluación médica.
Una función personalizada deberá diseñarse y validarse con profesionales
competentes y con los requisitos regulatorios aplicables. La versión inicial no
prometerá resultados de salud.

\subsection{Brechas para operación comercial}

Antes de atender clientes reales se requiere: desplegar o sustituir los workers
programados; disponer de Redis gestionado si se mantiene Celery; resolver
almacenamiento persistente de media; definir respaldos y recuperación; formalizar
soporte, términos y privacidad; ejecutar QA de navegador; y aislar los datos por
gimnasio. El último punto es crítico: la arquitectura actual no debe venderse como
multiempresa hasta implementar y probar separación de cuentas o tenants.
